\chapter{Subfactors, centers, and the tube algebra}
\label{ch:centers}

The six Haagerup lines are not braided, yet the chiral algebra in the
construction has a braided category of twelve representations.  This
chapter explains the bridge.  The main idea is that a charge must
specify how it passes \emph{every} symmetry line.  This extra data is
a half-braiding, and the objects carrying it form the Drinfeld center.
An equivalent description uses an algebra of annular line networks,
the tube algebra.  Before treating Haagerup, we work out the complete
example of $\Z_3$ symmetry, where all of these constructions reduce
to finite Fourier transforms.

\section{Why operator algebras enter the story}
\index{subfactor}\index{von Neumann algebra}\index{principal graph}

A von Neumann algebra is a unital adjoint-closed algebra of bounded
operators on a Hilbert space, closed in the weak operator topology:
convergence of every matrix element. It is a factor when its algebraic
center consists only of scalar multiples of the identity. Such algebras
arise naturally as local observable algebras in quantum field theory.
A subfactor is an inclusion $N\subset M$ of factors.  Finite index
measures the size of $M$ as an $N$-module, with the ordinary dimension
replaced by a dimension compatible with the operator-algebra trace or
standard representation.  For an outer action of a finite group $G$,
the fixed-point inclusion $M^G\subset M$ has index $|G|$.  This
recovers an integer in the group case, while general subfactors can
have nonintegral index.
Here an outer group action means that no nonidentity group element
acts as $x\mapsto uxu^*$ for a unitary $u$ inside $M$.

One should not try to obtain nonintegral index by counting an ordinary
basis of $M$ over $N$.  The relevant module dimension is not the
dimension of a finite-dimensional complex vector space.  The same
distinction already appeared when the positive dimension of a
Fibonacci line was $\varphi$.

An inclusion has a basic bimodule, conventionally written
${}_N L^2(M)_M$.  Fusing bimodules over the intermediate algebra
produces more bimodules.  Alternating this basic bimodule with its
dual generates tensor categories of $N$--$N$ and $M$--$M$
bimodules, called the even parts.  In a finite-depth situation there
are finitely many simple types, giving fusion categories.  Positivity
comes from the Hilbert-space bimodule structure.

Write $X={}_N L^2(M)_M$ for the basic bimodule. The
\emph{principal graph} records the decomposition process. Its even
vertices are the simple $N$--$N$ bimodules and its odd vertices the
simple $N$--$M$ bimodules generated by alternating $X$ with its dual.
An even vertex $Z$ is joined to an odd vertex $Y$ with multiplicity
$\dim\Hom(Y,Z\otimes_N X)$. The root is $L^2(N)$; reversing $N$
and $M$ gives the dual principal graph. Finite depth means that these
generated graphs have finitely many vertices.

The original Haagerup subfactor has index
$(5+\sqrt{13})/2$.  Its construction produces categorical data
that are not explained by a finite group action.  The category
$\HH_2$ and its Morita partners, including $\HH_3$, arise in this
setting.  For our purposes the key lesson is a method: a concrete
operator-algebra inclusion can package an exotic fusion category in
the same way a group fixed-point inclusion packages ordinary
symmetry.

\section{A Jones projection and the origin of a loop parameter}

A short calculation makes the relation between index and diagrams
more concrete.  Work with a finite-index inclusion of type
$\mathrm{II}_1$ factors, which carry compatible normalized finite
traces. Here $L^2(M)$ is the completion of $M$ for
$\langle x,y\rangle=\operatorname{tr}(x^*y)$, and left and right
multiplication give its bimodule actions. There is a trace-preserving
conditional expectation $E_N:M\to N$.  It is an $N$-bimodule map:
$E_N(n_1xn_2)=n_1E_N(x)n_2$.  On $L^2(M)$ the map extending
$x\mapsto E_N(x)$ is the orthogonal projection $e_N$ onto
$L^2(N)$.

For $x\in M$, viewed as left multiplication, apply the composite
$e_Nxe_N$ to a vector $y\in M$.  The result is
\[
 E_N(xE_N(y))=E_N(x)E_N(y).
\]
Since $M$ is dense in $L^2(M)$, this proves the operator relation
\begin{equation}
 e_Nxe_N=E_N(x)e_N.
 \label{eq:Jones-compression}
\end{equation}
The algebra generated by $M$ and $e_N$ is the first basic
construction.  Iterating it produces a tower and Jones projections
$e_1,e_2,\ldots$.  For index $I=[M:N]$, their trace and local
relations are
\begin{equation}
 \operatorname{tr}(e_i)=I^{-1},\qquad
 e_i^2=e_i=e_i^\dagger,\qquad
 e_ie_{i\pm1}e_i=I^{-1}e_i,
 \label{eq:Jones-projections}
\end{equation}
while distant projections commute.  These relations express the
effect of creating a cup--cap pair, joining it to its neighbor,
and simplifying the resulting contractible loop.

Set $\delta_{\rm TL}=\sqrt I$ and
$U_i=\delta_{\rm TL}e_i$.  Then
\begin{equation}
 U_i^2=\delta_{\rm TL}U_i,\qquad
 U_iU_{i\pm1}U_i=U_i.
 \label{eq:TL-normalization}
\end{equation}
These are the Temperley--Lieb relations.  Closing a contractible
strand contributes the scalar $\delta_{\rm TL}$.  The loop
parameter is the square root of the subfactor index, not the
index itself.

\begin{example}[The first even object at the Haagerup index]
For the original Haagerup inclusion,
$I=1+d$, so $\delta_{\rm TL}=\sqrt{1+d}$.  The fundamental
shaded strand followed by its oppositely shaded dual has a unit
summand.  This composition returns to an even part: an $N$--$M$
bimodule composes with an $M$--$N$ bimodule.  Removing
that summand with the two-strand Jones--Wenzl projection
$1-e_1=1-\delta_{\rm TL}^{-1}U_1$ leaves categorical dimension
\[
 \delta_{\rm TL}^2-1=I-1=d.
\]
In the relevant even part this is the object $\rho$.  Thus the
two numbers $\sqrt{1+d}$ and $d$ have different diagrammatic
origins: the former belongs to a shaded fundamental strand, the
latter to the noninvertible even object.  The Temperley--Lieb
relations capture a useful subalgebra of diagram operations;
they do not supply every Haagerup junction or its full associator.
\end{example}

\section{Braiding, twists, and nondegeneracy}
\index{braiding}\index{hexagon identity}\index{ribbon twist}

A braiding is a natural family of isomorphisms
$c_{x,y}:x\otimes y\to y\otimes x$. Naturality says that exchanging
two lines commutes with applying morphisms to them. The two
\emph{hexagon identities} say that exchanging one line with a fused
pair agrees with two successive exchanges. Suppressing associators,
they read
\begin{align*}
 c_{x,y\otimes z}&=(1_y\otimes c_{x,z})(c_{x,y}\otimes1_z),\\
 c_{x\otimes y,z}&=(c_{x,z}\otimes1_y)(1_x\otimes c_{y,z}).
\end{align*}
With parentheses restored, the associators turn each equality into
a six-sided diagram. Thus the hexagons constrain braiding and
associativity together. A commutative fusion table is necessary for
a braiding, but does not supply these maps or identities.

In splitting bases $t_{ab}^{z,\mu}:z\to a\otimes b$, the exchange
has a matrix called an $R$-symbol:
\[
 c_{a,b}t_{ab}^{z,\mu}
 =\sum_\nu [R^{ab}_z]_{\nu\mu}t_{ba}^{z,\nu}.
\]
An $F$-matrix changes a fusion-tree basis; an $R$-matrix exchanges
two neighboring factors. They are different operations, even when
both happen to be scalars on a particular channel.

A ribbon structure includes compatible twists $\theta_x:x\to x$,
dualities, and traces. For simple $x$ the twist is a scalar. Its
balancing identity is
\[
 \theta_{x\otimes y}
 =c_{y,x}c_{x,y}(\theta_x\otimes\theta_y).
\]
Consequently double exchange on the simple output channel $z$ is
$\theta_z/(\theta_x\theta_y)$ times the identity. Double exchange
is called \emph{monodromy}; it need not be trivial even when the
objects being exchanged are identical.

An object is \emph{transparent} when its monodromy with every object
is the identity. A ribbon fusion category is modular when the only
transparent objects are direct sums of the unit; equivalently here,
its normalized Hopf-link matrix is invertible. This is the
nondegeneracy condition behind the modular matrices of the next
chapter. A symmetric category, such as ordinary finite-group
representations with the usual flip, instead makes every object
transparent. Apart from the trivial category, it is not modular.

Topological lines in a two-dimensional spacetime do not automatically
carry these exchange maps: moving intersecting lines through each
other is extra data. The noncommutative Haagerup fusion ring forbids
such a braiding on its six simple lines. The center construction
below supplies a different category whose objects carry precisely
the additional passing data.

\section{An algebra object remembers how an extension multiplies}

Suppose a larger chiral algebra $H$ contains a smaller one $U$.
Viewed as a $U$-module, $H$ is a sum of irreducible modules.  The sum
alone does not specify $H$: its vertex operators also multiply these
summands.  Categorically, this multiplication is an algebra object.

\begin{definition}
An algebra object $A$ in a monoidal category consists of maps
\[
 \mu:A\otimes A\longrightarrow A,
 \qquad \eta:\one\longrightarrow A
\]
satisfying associativity, with the category's associator inserted,
and the two unit axioms.  It is connected if
$\dim\Hom(\one,A)=1$.
\end{definition}

The associativity equation is
\[
 \mu(\mu\otimes1_A)
 =\mu(1_A\otimes\mu)\mathfrak a_{A,A,A}.
\]
Thus the existence of $A$ depends on the $F$-symbols.  The fact that
the fusion product of an underlying object contains the right
summands does not construct a multiplication satisfying this equation.

In a braided category, commutativity means
$\mu c_{A,A}=\mu$.  This is the categorical version of locality
for a chiral extension.  Separability means that multiplication has
a splitting compatible with the left and right $A$ actions.  It is
the property that keeps module categories semisimple in this setting.
In a unitary category the corresponding positive structure is often
called a \emph{Q-system}: a compatible dagger Frobenius algebra,\index{Q-system}
with a specialness normalization such as
$\mu\mu^\dagger=\beta 1_A$ for a positive scalar $\beta$.
The Frobenius condition says that the comultiplication
$\Delta=\mu^\dagger:A\to A\otimes A$ respects both $A$ actions.
Suppressing associators, it reads
\[
 (\mu\otimes1_A)(1_A\otimes\Delta)
 =\Delta\mu
 =(1_A\otimes\mu)(\Delta\otimes1_A),
\]
an equality of maps $A\otimes A\to A\otimes A$.
Thus $\beta^{-1}\Delta$ is a bimodule splitting of multiplication.
Normalization conventions vary, but the positivity and algebra
compatibility are substantive.

\begin{example}[The group algebra object]
In the category $\operatorname{Vec}_G$ of finite-dimensional
$G$-graded vector spaces with trivial associator, the object
$A=\bigoplus_{g\in G}g$ has a multiplication taking the
$g\otimes h$ summand to $gh$.  Associativity is exactly the group
law.  If the associator is instead twisted by a three-cocycle
$\omega(g,h,k)$, scalar multiplication coefficients must satisfy
an equation whose coboundary cancels $\omega$.  The same direct sum
may therefore fail to admit the untwisted multiplication.  An
associator can obstruct an extension even when the fusion rules look
unchanged.
\end{example}

\begin{example}[The smallest associator obstruction to gauging]
\label{ex:centers-Z2-obstruction}
\index{associator!three-cocycle}\index{gauging obstruction}
Take the two labels $0,1\in\Z_2$, with fusion given by addition
modulo two; write $g$ for the nonidentity simple object. All
fusion spaces are one dimensional. Define the associator scalar by
\[
 \omega(a,b,c)=(-1)^{abc},\qquad a,b,c\in\{0,1\}.
\]
Thus its only nontrivial value is $\omega(1,1,1)=-1$.
The pointed pentagon becomes the scalar identity
\[
 \omega(b,c,d)\omega(a,b+c,d)\omega(a,b,c)
 =\omega(a,b,c+d)\omega(a+b,c,d).
\]
On the left the exponent is $bcd+abd+acd+abc$; on the
right it is $abc+abd+acd+bcd$. They agree modulo two.
The unit constraints hold because the phase is one whenever
any argument is zero. We have therefore constructed a coherent
unitary associator with the same fusion rule $g^2=\one$ as
the untwisted pointed category.

The sign is not removed by choosing different normalized binary
junctions. If their phases are $u(a,b)$, with
$u(0,a)=u(a,0)=1$, the gauge factor at $(1,1,1)$ is
\[
 \frac{u(1,1)u(0,1)}{u(1,1)u(1,0)}=1.
\]
The pentagon solutions called three-cocycles are identified under
just these changes of junction basis; group cohomology packages
that distinction. This example exhibits it without requiring the
general cohomology machinery.

Now try to form $A=\one\oplus g$, with a unit-normalized
multiplication and coefficient $\beta$ on $g\otimes g\to\one$.
Associativity restricted to $g\otimes g\otimes g$ says
\[
 \beta=-\beta.
\]
The left-associated product contributes $\beta$, while the
right-associated product includes the associator $-1$.
Consequently $\beta=0$ over $\C$. The resulting algebra has
the square-zero ideal $g$ and is not separable. Explicitly, its
right regular module has quotient $\one$ on which $g$ acts
by zero. A splitting would have to embed that quotient as the
unit summand of $A$, but right multiplication by $g$ takes
that summand to the nonzero $g$ summand. Thus the quotient
does not split as an $A$-module.

There is therefore no separable group-algebra condensation on
this sum. When the two objects represent symmetry defects, this
is an associator obstruction to gauging the whole $\Z_2$.
No braiding was required for the obstruction. It concerns an
algebra in the symmetry fusion category; a bosonic chiral
extension in a braided representation category also requires
the separate commutativity and locality conditions. For the
Haagerup construction, the actual order-three subgroup is proved
untwisted in Chapter~\ref{ch:branching}; its group law alone
would not have proved that fact.
\end{example}

For the original Haagerup subfactor, the connected algebra object
$\one\oplus\rho$ in the relevant $\HH_2$ category has
dimension $1+d$, agreeing with the subfactor index.  Its multiplication
is genuine additional data.  The different associator of $\HH_3$
prevents transferring that particular statement merely by copying
the two summand labels.

\section{Modules, local modules, and bimodules answer different questions}
\index{local module}

A right $A$-module is an object $X$ with an action
$r_X:X\otimes A\to X$ obeying the usual associative and unital
identities.  In a braided category, a commutative separable algebra
allows relative tensor products over $A$, giving a category of all
right modules.  The braiding moves an $A$ line to the side where it
can act; one then identifies the two ways of applying that action.
This is the categorical analogue of balancing a tensor product over
an ordinary algebra.
Our convention equips a right module with the left action
$\lambda_X=r_Xc_{X,A}^{-1}:A\otimes X\to X$ when forming
relative tensor products; this fixes the braid direction in
Theorem~\ref{thm:lagrangian-reconstruction}.

Not every module is \emph{local}.  A right module is local, or
dyslectic, when moving $A$ all the way around it before applying
the action has no effect:
\begin{equation}
 r_X c_{A,X}c_{X,A}=r_X.
 \label{eq:local-module-condition}
\end{equation}
The order of maps starts at $X\otimes A$.  Local modules inherit
a braiding and describe ordinary representations of a local chiral
extension.  Nonlocal modules still exist and can describe boundary
or defect data, even when the extended chiral algebra is holomorphic.

There are therefore three related categories:
\begin{center}
\begin{tabular}{c|p{0.62\textwidth}}
Category & Role in an extension\\\hline
All right $A$-modules & Boundary or one-sided extension data;
may be a nonbraided fusion category.\\
Local right $A$-modules & Representations of the extended local
chiral algebra.\\
$A$--$A$ bimodules & Defects compatible with the chosen extension
description, with composition by relative tensor product.
\end{tabular}
\end{center}
For the selected holomorphic extension in this book, the local
module category is $\operatorname{Vec}$, while the category of all
right modules is $\HH_3$.  The word ``holomorphic'' therefore does
not mean that the extended theory has no interesting defects.

\section{Morita equivalence changes the boundary description}
\index{Morita equivalence}

For ordinary algebras, $M_n(\C)$ and $\C$ have equivalent module
categories: an $M_n(\C)$-module is a sum of copies of the standard
column module, and its multiplicity space is a vector space.  The
algebras are not isomorphic, but they present the same module theory.
This is the elementary model of Morita equivalence.

\index{module category}\index{dual category}
For fusion categories the analogue of a representation is a
\emph{module category}. In a right-action convention it consists
of a finite semisimple category $\mathcal N$, an action
$(m,X)\mapsto m\triangleleft X$, and coherent isomorphisms
\[
 (m\triangleleft X)\triangleleft Y
   \simeq m\triangleleft(X\otimes Y),\qquad
 m\triangleleft\one\simeq m.
\]
Their coherence is the module version of the pentagon. A module
endofunctor $T$ also has coherent isomorphisms
$T(m\triangleleft X)\simeq T(m)\triangleleft X$:
it respects the action, rather than merely mapping objects to objects.
Such functors form the \emph{dual category}, with composition as
tensor product. If $\mathcal N$ is indecomposable as a module
category, this dual is a fusion category. For left module categories,
the corresponding dual-category convention uses reverse composition.

For a small example, let untwisted $\operatorname{Vec}_{\Z_3}$
act on $\operatorname{Vec}$ by forgetting the grading. Each simple
group label acts as the identity functor. A linear endofunctor is
$T_W(v)=W\otimes v$; respecting the generator's action equips
$W$ with an operator $U$ satisfying $U^3=1$. The three one-dimensional
eigenvalue choices $1,\zeta_3,\zeta_3^2$, where
$\zeta_3=e^{2\pi i/3}$, are the three simple module
endofunctors, and composition multiplies their characters. Thus the
module category has one simple object, whereas its dual has three.
Fusion categories related through such an indecomposable module
category are Morita equivalent. Their simple objects and fusion
rings can differ; their centers agree by the theorem below.

\begin{theorem}[Morita invariance of the center]
Morita equivalent fusion categories have braided equivalent
Drinfeld centers.  In the unitary setting the equivalence can be
taken compatible with the unitary structures and canonical ribbon
structures.
\end{theorem}

We use this structural theorem as an input from tensor-category
theory.  Its physical interpretation is helpful: distinct boundary
descriptions of the same three-dimensional topological theory can
have the same bulk anyons.  Those bulk anyons form the center.
Thus $Z(\HH_2)\simeq Z(\HH_3)$ is entirely compatible with
the inequivalence demonstrated by \eqref{eq:H2-H3-invariant}.

Morita equivalence does not itself give an isospectral map between
arbitrarily chosen lattice Hamiltonians.  To compare chains one
would need a specified construction transporting their Hilbert
spaces, local interactions, boundary conditions, and inner products.
Center equivalence transports categorical bulk data, not an
unspecified dynamical model.

\section{A center object knows how to pass every line}
\index{Drinfeld center}\index{half-braiding}

\begin{definition}
An object of the Drinfeld center $Z(\mathcal C)$ is a pair
$(z,\sigma)$ consisting of an object $z$ of $\mathcal C$ and
a natural family of isomorphisms
\[
 \sigma_x:z\otimes x\longrightarrow x\otimes z
 \qquad(x\in\mathcal C),
\]
called a half-braiding.  It is compatible with tensor products:
passing $z$ through $x\otimes y$ agrees with passing it through
$x$ and then through $y$, with the associators inserted.
\end{definition}

In a strict notation the compatibility equation reads
\begin{equation}
 \sigma_{x\otimes y}
 =(1_x\otimes\sigma_y)(\sigma_x\otimes1_y).
 \label{eq:half-braiding-strict}
\end{equation}
Strict notation suppresses associators; it does not remove them
from a calculation with actual $F$-matrices.  Naturality means that
the same motion commutes with every junction $x\to y$.

Two center objects can be braided: use the first object's
half-braiding with the underlying object of the second.  This makes
$Z(\mathcal C)$ braided even when $\mathcal C$ is not.  With a
spherical structure in the usual characteristic-zero setting, the
center is modular.  Its global dimension is
\begin{equation}
 \operatorname{Dim}(Z(\mathcal C))
 =\operatorname{Dim}(\mathcal C)^2.
 \label{eq:center-global-dimension}
\end{equation}
The forgetful functor $F:Z(\mathcal C)\to\mathcal C$ discards
the half-braiding.  A simple center object can become a nonsimple
object after forgetting; this is precisely what happens for most
Haagerup center objects.

\section{The complete \texorpdfstring{$\Z_3$}{Z3} example}

Take $\mathcal C=\operatorname{Vec}_{\Z_3}$ with trivial
associator.  Its simple objects are $g\in\Z_3$, with fusion
$g\otimes h=g+h$.  Because the group is abelian, a half-braiding
on simple object $g$ with simple object $h$ is a scalar
$\sigma_h$.  Equation~\eqref{eq:half-braiding-strict} says
$\sigma_{h+h'}=\sigma_h\sigma_{h'}$, so it is a character:
\[
 \sigma_h=\omega^{kh},\qquad k\in\Z_3.
\]
Here $\omega=e^{2\pi i/3}$. On an arbitrary graded object, the
half-braiding with the generator acts on each graded multiplicity
space by an operator of order three. It diagonalizes into the three
characters, so every center object decomposes into these pairs $(g,k)$.
There are therefore nine simple center objects, labelled $(g,k)$.
The label $g$ is a flux or twisted-sector label, and $k$ is a charge.
Their fusion and dimensions are
\begin{equation}
 (g,k)\otimes(h,l)=(g+h,k+l),\qquad d_{(g,k)}=1.
 \label{eq:Z3-center-fusion}
\end{equation}
The twist, obtained by closing the half-braiding of the object with
itself, is
\begin{equation}
 \theta_{(g,k)}=\omega^{gk}.
 \label{eq:Z3-center-twist}
\end{equation}

The braiding is just as explicit. Passing $(g,k)$ through $(h,l)$
uses the first object's character on the second object's underlying
flux, so its scalar is $\omega^{kh}$. Passing back gives
$\omega^{lg}$, hence the monodromy is $\omega^{kh+lg}$.
This agrees with balancing, since
\[
 \frac{\theta_{(g+h,k+l)}}{\theta_{(g,k)}\theta_{(h,l)}}
 =\omega^{(g+h)(k+l)-gk-hl}=\omega^{kh+lg}.
\]
It also proves nondegeneracy directly. A transparent $(g,k)$ must
have trivial monodromy with $(1,0)$, which forces $k=0$, and with
$(0,1)$, which forces $g=0$. Only the unit is transparent.

To fix a torus convention that we will maintain in the next chapter,
take
\begin{equation}
 S_{(g,k),(h,l)}=\frac13\omega^{-(kh+lg)}.
 \label{eq:Z3-center-S}
\end{equation}
The minus sign corresponds to the inverse-monodromy orientation of
the Hopf link and is compatible with the displayed positive twist.
Using the opposite Hopf-link orientation requires the corresponding
change in the modular generators.  On the real, self-dual Haagerup
data this orientation issue is invisible, which makes the present
example a useful convention check.

Finite Fourier orthogonality immediately gives
\[
 S^2_{(g,k),(h,l)}=\delta_{h,-g}\delta_{l,-k}.
\]
Thus $S^2$ is charge conjugation rather than the identity.  Also
\[
 \sum_{g,k}\theta_{(g,k)}=\sum_g\sum_k\omega^{gk}=3.
\]
Dividing by total quantum dimension $3$ gives Gauss--Milgram phase
one.  Direct finite Fourier multiplication gives $(S\Theta)^3=S^2$,
where $\Theta$ is the twist matrix.  These statements will be
generalized in Chapter~\ref{ch:modular}.

\section{The tube algebra in the same example}
\index{tube algebra}

An annular diagram can carry an incoming defect label $a$, an
outgoing label $b$, and a line $c$ circling the annulus.  Its junction
belongs to
\[
 \Hom(c\otimes a,b\otimes c).
\]
Taking the direct sum over simple labels gives the vector space
\begin{equation}
 \operatorname{Tube}(\mathcal C)
 =\bigoplus_{a,b,c}\Hom(c\otimes a,b\otimes c).
 \label{eq:tube-vector-space}
\end{equation}
Multiplication glues annuli.  Resolving the two circling lines into
simple channels uses fusion and the associator.  This multiplication
is the part that an ordinary list of fusion coefficients does not
specify.  In a unitary category, reflection gives a positive adjoint
structure, so the tube algebra is a finite-dimensional $C^*$-algebra.

For $\operatorname{Vec}_{\Z_3}$, the morphism space is nonzero
exactly when $a=b$.  Let $e_{g,h}$ denote the annulus with defect
label $g$ and circling line $h$.  Then
\begin{equation}
 e_{g,h}e_{g',h'}=\delta_{g,g'}e_{g,h+h'}.
 \label{eq:Z3-tube-product}
\end{equation}
Thus the tube algebra is three copies of $\C[\Z_3]$, one for
each flux.  Its nine primitive central idempotents are
\begin{equation}
 p_{g,k}=\frac13\sum_{h=0}^2\omega^{-kh}e_{g,h}.
 \label{eq:Z3-tube-projector}
\end{equation}
The identity
$\sum_h\omega^{(k-l)h}=3\delta_{kl}$ proves
$p_{g,k}p_{g',l}=\delta_{gg'}\delta_{kl}p_{g,k}$.
On the corresponding one-dimensional representation,
$e_{g,h}$ acts by $\omega^{kh}$.  A Dehn twist is represented by
$e_{g,g}$ and has eigenvalue $\omega^{kg}$, reproducing
\eqref{eq:Z3-center-twist}.

This calculation exhibits the equivalence between simple tube
representations and simple center objects without any large matrices.
The ordinary untwisted sector has $g=0$, so only three of the nine
labels occur there.  Looking only at operators $U_h$ on that sector
would reveal the three ordinary group characters, while omitting the
six charges attached to nontrivial fluxes.

\section{Changing the associator changes the center}

\begin{example}[Toric-code and double-semion twists]
\label{ex:centers-double-semion}
\index{double-semion category}\index{toric code}
Return to the two-object category of
Example~\ref{ex:centers-Z2-obstruction}, with
$\omega(a,b,c)=(-1)^{abc}$. All labels and sums below are in
$\Z_2$. For a center object with underlying label $g$, write
$\sigma_x$ for its scalar half-braiding with $x$.
Restoring the three associators in
\eqref{eq:half-braiding-strict} gives
\[
 \sigma_{x+y}
 =\frac{\omega(x,g,y)}
        {\omega(g,x,y)\omega(x,y,g)}\sigma_x\sigma_y
 =(-1)^{gxy}\sigma_x\sigma_y.
\]
Since $\sigma_0=1$, the choice $t=\sigma_1$ must satisfy
$t^2=(-1)^g$. There are two possibilities over each flux:
\[
 \one=(0,1),\qquad b=(0,-1),\qquad
 s=(1,i),\qquad \bar s=(1,-i).
\]
The corresponding quadratic operator on any multiplicity space
has distinct roots, so it diagonalizes into these choices.
They are the four simple center objects.

The tensor-product half-braiding also contains associators.
For $(g,t)$ and $(h,u)$ its value on the generator is
multiplied by
$\omega(g,h,1)\omega(1,g,h)/\omega(g,1,h)=(-1)^{gh}$.
Consequently
\[
 (g,t)\otimes(h,u)=(g+h,(-1)^{gh}tu),
\]
and hence
\[
 b^2=s^2=\bar s^{\,2}=\one,\qquad
 s\bar s=b,\qquad bs=\bar s,\qquad b\bar s=s.
\]
The fusion ring is that of $\Z_2\times\Z_2$.
In particular all four simples are self-dual. The bar on
$\bar s$ names the \emph{anti-semion}, whose twist is the
complex conjugate of the semion's; it does not denote the
categorical dual of $s$.

Use the canonical positive spherical structure. A direct sign
check is available in the left-dual convention of
Chapter~\ref{ch:fusion}: take
$\operatorname{coev}_g=1$,
$\operatorname{ev}_g=(-1)^g$, and pivotal map $j_g=(-1)^g$.
Their trace gives dimension one. The Drinfeld bend is the map
$g\to g^{**}$ made from a crossing, a cup, and a cap.
Its associator and evaluation signs cancel here, giving scalar
$u_g=\sigma_g$, so the ribbon twist is
$u_g^{-1}j_g=\sigma_g$, using $t^2=(-1)^g$.
Thus
\[
\begin{array}{c|rrrr}
 X&\one&b&s&\bar s\\\hline
 d_X&1&1&1&1\\
 \theta_X&1&1&i&-i
\end{array}
\]
These are the double-semion twists. Balancing makes the
monodromy of $b$ with either semion equal to $-1$.
Every nonidentity simple has a nontrivial monodromy with some
object, so the category is nondegenerate.

For comparison, the untwisted center
$Z(\operatorname{Vec}_{\Z_2})$ has labels $(g,k)$ with
ordinary componentwise fusion and twist $(-1)^{gk}$.
Its simples $1,e,m,\epsilon$ therefore have twists
$1,1,1,-1$, the toric-code pattern. Its fusion ring is again
$\Z_2\times\Z_2$, but no relabeling identifies these twists
with $1,1,i,-i$. Both centers have global dimension four and
Gauss--Milgram phase one:
\[
 \frac{1+1+i-i}{2}
 =\frac{1+1+1-1}{2}=1.
\]
Fusion multiplicities, global dimension, and the chiral anomaly
therefore leave information that the associator and center twists
can distinguish. This is a small complete model of why actual
Haagerup junction data are needed.
\end{example}

\section{Haagerup forgetful multiplicities}

The tube-sector decomposition calculated for $\Z_3$ also occurs
for Haagerup symmetry, with a less trivial distribution of
multiplicities. Use center labels
\[
 0,b,a,c_0,c_1,c_2,d_1,\ldots,d_6.
\]
Here $a$ is a center object; it is not the invertible object
$\alpha$ of $\HH_3$.  Put $X=\rho\oplus\alpha\rho\oplus
\alpha^2\rho$.  The forgetful images are
\begin{align}
 F(0)&=\one,&F(b)&=\one\oplus X,&F(a)&=2\one\oplus X,\notag\\
 F(c_j)&=\alpha\oplus\alpha^2\oplus X,&
 F(d_j)&=X.
 \label{eq:H3-forgetful}
\end{align}
These are categorical center data.  Their dimensions already give
the complete list of central quantum dimensions:
\[
 1,\quad1+3d=d^2,\quad2+3d=s\quad
 \hbox{(four objects)},\quad3d\quad\hbox{(six objects)}.
\]

For a simple defect $x$ define the finite-dimensional space
\begin{equation}
 W_x^\mu=\Hom_{\HH_3}(F(\mu),x),
 \qquad n_{x\mu}=\dim W_x^\mu.
 \label{eq:H3-W-definition}
\end{equation}
The multiplicity table, with identical columns grouped, is
\begin{center}
\begin{tabular}{c|ccccc}
$x$ & $0$ & $b$ & $a$ & each $c_j$ & each $d_j$\\\hline
$\one$ &1&1&2&0&0\\
$\alpha$ &0&0&0&1&0\\
$\alpha^2$ &0&0&0&1&0\\
each $\alpha^g\rho$ &0&1&1&1&1
\end{tabular}
\end{center}
The irreducible tube module $W^\mu=\bigoplus_xW_x^\mu$ therefore
has ordinary vector-space dimension
\begin{equation}
 1,\ 4,\ 5,\ 5,\ 5,\ 5,\ 3,\ 3,\ 3,\ 3,\ 3,\ 3.
 \label{eq:H3-tube-dimensions}
\end{equation}
These integers are not the central quantum dimensions.

\begin{example}[The dimension of the Haagerup tube algebra]
Finite-dimensional $C^*$-algebra structure gives
\begin{equation}
 \operatorname{Tube}(\HH_3)
 \simeq\C\oplus M_4(\C)\oplus M_5(\C)^{\oplus4}
 \oplus M_3(\C)^{\oplus6}
 \label{eq:H3-tube-Wedderburn}
\end{equation}
as an algebra.  Its dimension is
\[
 1^2+4^2+4\cdot5^2+6\cdot3^2=171.
\]
This finite integer should be compared with neither $3s$ nor
$(3s)^2$.  Those are categorical global dimensions.  The tube algebra
is an ordinary algebra built from junction spaces and has an ordinary
vector-space dimension.
\end{example}

The center equivalence gives its representations additional tensor
and braiding structure.  An abstract algebra isomorphism such as
\eqref{eq:H3-tube-Wedderburn} by itself forgets that structure.
Just as knowing $M_n(\C)$ does not tell us which preferred matrices
represent a Hamiltonian, knowing the matrix-block sizes does not
reconstruct the categorical junction operations.

\section{Finite tube factors versus physical multiplicity spaces}

In a two-dimensional theory with $\mathcal C$ symmetry, let
$\mathscr H_x$ be the Hilbert space on a circle intersected by a
defect $x$.  Its decomposition takes the form
\begin{equation}
 \mathscr H_x\simeq\bigoplus_{\mu\in\operatorname{Irr}Z(\mathcal C)}
 W_x^\mu\otimes\mathcal V_\mu.
 \label{eq:physical-tube-decomposition}
\end{equation}
Here $\operatorname{Irr}Z(\mathcal C)$ denotes representatives of the
simple isomorphism classes. The space $W_x^\mu$ is finite and
categorical. The space
$\mathcal V_\mu$ contains the dynamical states: in a CFT it carries
energies, spins, and descendant towers.  The tube algebra acts on
the first factor, while the Hamiltonian commutes with that action
and acts on the multiplicity space.

For Haagerup symmetry, the ordinary untwisted space is
\begin{equation}
 \mathscr H_{\one}\simeq
 \mathcal V_0\oplus\mathcal V_b\oplus
 (\C^2\otimes\mathcal V_a).
 \label{eq:H3-untwisted-decomposition}
\end{equation}
Its three types agree with the two scalar representations and the
doublet of Chapter~\ref{ch:haagerup}.  The factor of two is a real
physical multiplicity in the full untwisted Hilbert space.  A
partition function defined for the multiplicity space
$\mathcal V_a$ alone must not include it a second time.

For the selected chiral Haagerup action on the holomorphic $E_8$
full theory, the construction will identify
\[
 \mathcal V_\mu=M_\mu\otimes\overline{V_{E_8}}.
\]
That is a dynamical realization of
\eqref{eq:physical-tube-decomposition}, not a consequence of the
abstract finite fusion ring by itself.

\section{Why the tube algebra is needed for operator products}

For an ordinary group, the tensor product of representations is
defined by $U_g\mapsto U_g\otimes U_g$.  Trying the same rule for
a noninvertible fusion algebra fails.  If $Y_\rho=d$ on two
one-dimensional invariant representations, the naive tensor-product
rule would give $Y_\rho=d^2$.  But $d^2$ is neither of the allowed
scalar roots of $y^2=1+3y$.  The putative tensor product is not even
a representation of the fusion algebra.

The missing ingredient is the junction information specifying how
lines act around two operators and how the two lassos are joined.
The center, or equivalently the tube algebra equipped with its
categorical structure, supplies that information.  Local operators
transform in center charges.  If the symmetry is unbroken, a
correlator with charges $\mu_1,\ldots,\mu_n$ can be nonzero only
if
\begin{equation}
 \Hom_{Z(\mathcal C)}
 (\one,\mu_1\otimes\cdots\otimes\mu_n)\ne0.
 \label{eq:center-selection-rule}
\end{equation}
This is a selection rule, not a formula for a normalized three-point
coefficient.  Allowed tensor channels still carry dynamical
amplitudes.  Extracting such amplitudes from a lattice calculation
requires identified operator families and coherent normalizations,
not only closed-line eigenvalues.

\section{Exercises}

\begin{exercise}\label{ex:Z3-halfbraid}
Starting from \eqref{eq:half-braiding-strict}, classify the
half-braidings on each simple object of
$\operatorname{Vec}_{\Z_3}$.  Compute all twists and count how
many are $1$, $\omega$, and $\omega^2$.

\par\smallskip\noindent{\footnotesize\normalfont
\hyperref[sol:ex:Z3-halfbraid]{Solution on page~\pageref*{sol:ex:Z3-halfbraid}.}}\par
\end{exercise}

\begin{exercise}\label{ex:Z3-tube}
Prove the projector multiplication rule following
\eqref{eq:Z3-tube-projector}.  Show that the sum of all nine
projectors is the unit of the tube algebra.  Which three projectors
can act on the untwisted sector?

\par\smallskip\noindent{\footnotesize\normalfont
\hyperref[sol:ex:Z3-tube]{Solution on page~\pageref*{sol:ex:Z3-tube}.}}\par
\end{exercise}

\begin{exercise}\label{ex:tube171}
Verify the dimension $171$ directly from
\eqref{eq:tube-vector-space} and the Haagerup fusion rules, using
\[
 \dim\Hom(ca,bc)=\sum_z N_{ca}^{\ z}N_{bc}^{\ z}.
\]
First sum over the three invertible choices of $c$, then over the
three noninvertible choices.

\par\smallskip\noindent{\footnotesize\normalfont
\hyperref[sol:ex:tube171]{Solution on page~\pageref*{sol:ex:tube171}.}}\par
\end{exercise}

\begin{exercise}\label{ex:forgetful-qdim}
Use \eqref{eq:H3-forgetful} to derive every central quantum
dimension.  Compare the quantum dimension and ordinary tube-module
dimension of the object $a$.  Why are these allowed to differ?

\par\smallskip\noindent{\footnotesize\normalfont
\hyperref[sol:ex:forgetful-qdim]{Solution on page~\pageref*{sol:ex:forgetful-qdim}.}}\par
\end{exercise}

\begin{exercise}\label{ex:all-local}
Explain how a holomorphic chiral extension can have a trivial
category of local modules and a nontrivial category of all modules.
Identify the extra equation a module must satisfy to be local, and
explain why a direct-sum decomposition of the extension is
insufficient to reconstruct its multiplication.

\par\smallskip\noindent{\footnotesize\normalfont
\hyperref[sol:ex:all-local]{Solution on page~\pageref*{sol:ex:all-local}.}}\par
\end{exercise}

\paragraph{Source note.}
The center--tube correspondence and the physical multiplicity-space
interpretation are standard structural results.  The treatment of
symmetry charges follows the formulation of Lin, Okada, Seifnashri,
and Tachikawa.  The Haagerup forgetful table is supplied as established
center data; the explicit decompositions and dimension calculations
following it are derived on the page.
